\section*{Introduction}

Bacterial antimicrobial resistance was associated with an estimated 4.95 million deaths in 2019
\cite{murray2022gram}, and the burden is forecast to grow to 2050 \cite{naghavi2024forecast}.
Whole-genome sequencing offers susceptibility information without waiting for culture, in one of
two forms. Rule-based tools look up catalogued determinants and call a strain resistant when one is
found \cite{bortolaia2020resfinder,feldgarden2019amrfinder,feldgarden2021amrfinderplus}; learned
models instead fit the genotype--phenotype map from labelled genomes
\cite{moradigaravand2018prediction,her2018pangenome,ren2022giessen,orcales2024tutorial}. Learned
models routinely report high accuracy. External validation of both approaches, however, finds error
rates that differ by drug and by collection \cite{verschuuren2022external}, and part of the
resistance in \Ecoli{} --- to amoxicillin--clavulanic acid above all --- is explained by no
catalogued variant \cite{lipworth2024uncatalogued}.

What makes those accuracy figures hard to read is population structure. Bacteria reproduce
clonally, so distant variants stay linked and whole lineages inherit their genomes and, often,
their resistance phenotypes together. Split such a collection at random and a model can score well
by recognising lineages rather than mechanisms. Yu, Wheeler and Barquist showed across five
pathogens that biased sampling driven by population structure confounds machine-learning
prediction of resistance, and that more data does not fix it \cite{yu2025biased}. Benchmarks that
compare learned models with rule-based ResFinder report that the ranking of the two depends on how
closely the test genomes are related to the training genomes \cite{hu2024assessing}, and
genotype--phenotype discordance is itself drug- and species-specific \cite{cheon2026discordance}.

Three questions follow that no single random-split benchmark can answer. First, how much of a
model's apparent accuracy comes from information that is not resistance biology at all: strain
metadata, study of origin, or properties of the assembly? Second, when lineages are held out, does
prediction fail because the model cannot generalise, or because the features it was given do not
contain the relevant determinants? Third, once the determinants are present, does a learned model
add anything over the obvious alternative, a fixed rule that calls a strain resistant when a known
determinant is detected?

Answering these requires a model good enough that its failures are informative rather than its own
fault, so we first built one. TriLens reads three complementary views of each genome --- mutation
patterns inside genes, gene presence and absence, and coarse per-gene statistics --- through three
neural branches joined by a learned gate, and stacks that with five classical learners under a
logistic-regression meta-learner. On a random split it leads every external baseline we tried. We
then used it, alongside simpler models and a determinant rule, to take the random split apart.

We work on the public collection of 6,962 \Ecoli{} genomes with phenotypes released by Arcadia
Science \cite{arcadia2024dataset,bell2025arcadia}, with an independent cohort of 10,232 isolates
from NCBI Pathogen Detection \cite{ncbi_pathogen_detection}, their minimum inhibitory
concentrations, a second independent \Ecoli{} collection \cite{ren2022giessen}, and a second
species. Every analysis bearing on a claim about features, rules or external transfer was
pre-registered: its endpoint and pass rule were committed to a public repository, with a timestamp,
before the relevant data were read, and every outcome is reported, including the four that failed.
Our contributions are:

\begin{enumerate}
  \item a three-view stacked model with a learned per-drug gate, and an ablation of each of its
        parts, reported against external baselines and against its own strongest component;
  \item a decomposition of random-split accuracy into metadata, assembly and genomic signal, with a
        test of whether the assembly signal is a batch artefact or biology;
  \item a prevalence correction for lineage-aware testing, reported per drug, which separates the
        part of a holdout's drop that is the metric from the part that is the model;
  \item a pre-registered demonstration that isolate-specific determinant calls rescue prediction on
        unseen lineages, with its coverage limits stated;
  \item a frozen external validation on a cohort from different countries, sources and years, with
        a leave-one-BioProject-out audit and MIC-based labels harmonised to EUCAST breakpoints;
  \item a comparison of learned models against a determinant rule in two species, showing that a
        sparse linear model on determinant calls suffices where a model helps at all, and
        identifying where it does not help;
  \item a diagnosis of a degenerate gated-attention score on presence/absence input, and negative
        results reported as such: the gate does not route by mechanism, attention recovers no
        catalogued resistance gene, and no approach meets clinical error benchmarks.
\end{enumerate}

\section*{Related work}

\subsection*{What models are fed, and how they are combined}
Studies of genomic resistance prediction differ mainly in the representation they feed a model
(Table~\ref{tab:summary}): pan-genome gene presence and absence
\cite{moradigaravand2018prediction,her2018pangenome}, $k$-mers \cite{jain2025paeruginosa,
gao2024abaumannii,serajian2024mtb}, unitigs \cite{do2024unitig}, SNPs and variant calls
\cite{wang2025hpylori,paredes2025mtb}, or mass spectra \cite{lopezcortes2024msdeepamr}. Across them,
gradient boosting and random forests remain hard to beat, which matches what tabular benchmarks
outside this field report \cite{grinsztajn2022trees,shwartzziv2022tabular}; TabNet was designed to
close that gap with attention over tabular inputs \cite{arik2021tabnet}, and we use it as a
baseline for exactly that reason.

Our design differs in how it combines representations rather than in which one it uses, so the
methods it draws on are worth naming. The mixture-of-experts idea --- several specialised
sub-models and a gate that learns how much to trust each --- goes back to Jacobs et al.\
\cite{jacobs1991moe} and was scaled to modern networks by Shazeer et al.\ \cite{shazeer2017moe}.
Its appeal here is direct: different drugs are resisted through different mechanisms, so a fixed
combination of views dilutes whichever view is right for a given drug. Wolpert's stacked
generalisation uses base-learner predictions as features for a second-level model
\cite{wolpert1992stacked}; fitting that level on out-of-fold predictions keeps it from trusting a
learner that has memorised its training data. Gated attention pooling \cite{ilse2018attention}, a
variant of attention \cite{vaswani2017attention} built for bags in which few items are informative,
describes a pan-genome well, where a strain carries thousands of genes and a handful matter. For
interpretation, SHAP provides attribution with consistency guarantees \cite{lundberg2017shap}, and
we use it alongside the model's own attention weights rather than in place of them.

Every one of those methods was developed where training and test samples are drawn independently
from one distribution. That is the assumption under which a gate learns which sub-model to trust
and a meta-learner learns how to weight its inputs. Whether a learned gate still routes sensibly
when the test set comes from an unseen lineage, and whether the neural and classical halves of a
stack degrade at the same rate, are not questions this literature has asked. We answer both for our
own model, and the answers are not the ones its architecture assumed.

\subsection*{Evaluation design and population structure}
The problem our evaluation addresses is old and well understood in bacterial genome-wide
association studies, a literature largely separate from predictive machine learning. Saber and
Shapiro benchmarked association methods on simulated bacterial genomes and phenotypes precisely
because population structure and genome-wide linkage make a causal variant hard to separate from
one merely linked to it \cite{saber2020bacgwas}. The response there has been to build the
correction into the model: Lees et al.\ fit an interpretable pangenome-spanning regression in which
variants compete inside one model instead of being tested one at a time, absorbing much of the
signal that would otherwise be attributed to a linked variant \cite{lees2020pyseer}.

Predictive work has largely not adopted that correction, and the reason is structural rather than
careless. An association study asks which variant causes the phenotype, so confounding by ancestry
is its central threat. A prediction study asks how well the phenotype can be guessed, and there
ancestry is a legitimate source of predictive signal, so nothing in the task itself forces the
question. The confound becomes visible only when a model meets a lineage it has not seen, and a
random split never produces that case.

The evaluation-split column of Table~\ref{tab:summary} shows how that plays out. Fourteen of the
sixteen entries predict a strain's phenotype; the other two are a review and a tool that classifies
genes into resistance categories. Of the fourteen, two partition by population structure as their
primary design \cite{yu2025biased,hu2024assessing}, and both report a drop once they do. Five pair
cross-validation or a holdout with an independent external collection, three use cross-validation
alone, and four use a random or stratified holdout with no external set. Two report a
structure-aware analysis beside a main result that does not use one: Moradigaravand et al.\ hold out
ST131 in a supplementary analysis, where the holdout scores about 0.28 below a size-matched random
control \cite{moradigaravand2018prediction}, and Orcales et al.\ offer MLST-blocked
cross-validation as an option not used for the headline figure \cite{orcales2024tutorial}. This is
not a criticism of any individual study, and our own first attempt at a structure-aware split had
the same defect. But none of the fourteen corrects the measure for the prevalence shift a lineage
holdout causes, and none separates the resulting drop by antibiotic.

Two gaps remain even after Yu et al. Their result is an average over five pathogens, so it
establishes that a drop exists without saying which drugs carry it. And accounting for structure
while splitting is not the same as holding a lineage out: a split can balance lineage composition
across partitions and still leave close relatives on both sides, which measures something else
entirely. Both gaps are what this paper's prevalence-corrected, per-drug holdout is for.
